\section{Coherent clusters of work heads}\label{sec:clusters}

Several physical heads can share the protected-slab argument. Let
$\varnothing\ne F\subseteq\{1,\ldots,k\}$, $m=|F|$, choose fixed signs
$\sigma_h\in\{-1,1\}$, and put $u_h(t)=\sigma_hy_h(t)$ and
$z(t)=\sum_{h\in F}u_h(t)$. A normalized chart consists of nondecreasing
$1$-Lipschitz maps $f_h:\R\to\R$ with $\sum_h f_h(v)=v$ for every real
$v$. Its error is at most $J$ if $|u_h(t)-f_h(z(t))|\le J$ at every live
step. The cluster, signs and maps may depend on the input, but remain fixed
through that execution.

\begin{lemma}[Protected cores and paid fringes]\label{lem:cores}
\formal{ProtectedCores.lean}{selected_in_core_or_fringe,outside_avoids_core,protected_core_disjoint_fringe,normalized_chart_core_fringe_counts}
For $J\ge0$, let $j=\lceil J\rceil$ and select a carrier slab $[\ell,\ell+W)$, and put
$A_h=f_h(\ell)$, $B_h=f_h(\ell+W)$. On oriented tape $h$ define
\[
 C_h=(A_h+j,B_h-j)\cap\Z,\qquad
 F_h=([A_h-j,A_h+j]\cup[B_h-j,B_h+j])\cap\Z.
\]
Every selected read lies in $C_h\cup F_h$; every outside-slab step avoids
$C_h$; $C_h\cap F_h=\varnothing$; and
$\sum_h|C_h|\le W+m$, $\sum_h|F_h|\le m(4j+2)$.
Only four rounded endpoints per head are needed to encode these sets.
\end{lemma}
\begin{outline}
Inside the carrier slab monotonicity places $f_h(z)$ between $A_h$ and
$B_h$, and the error places the physical position in $[A_h-j,B_h+j]$.
Its complement to the strict core lies in the two closed fringes, including
when they overlap or the core is empty. If $z<\ell$, the position is at
most $A_h+j$; if $z\ge\ell+W$, it is at least $B_h-j$. Thus no outside
old position writes a core. A real interval of length $B_h-A_h$ contains
at most that length plus one integer. Sum and use
$\sum_h(B_h-A_h)=W$. Each closed fringe interval contains at most $2j+1$
integers. Rounding by floor and ceiling determines all endpoints exactly;
reflection by $\sigma_h$ preserves the counts and ownership. Equality at a
core boundary is deliberately assigned to the fringe. The checked bound also
allows individual errors $j_h$, with total fringe bill $\sum_h(4j_h+2)$.
\end{outline}

\begin{theorem}[Coherent-cluster extension]\label{thm:clusters}
\formal{ClusterExpansion.lean}{native_coherent_cluster_returns,native_coherent_cluster_expansion}
The conclusions of \cref{thm:returns,thm:off-carrier} extend to every
$\varepsilon P$-coherent cluster, with returns and target lower endpoints on
its complement. In the return statement one may take
\[
 N_c=\lfloor160akC^2\rfloor+1,\quad
 \varepsilon=(128akN_c^2)^{-1},\quad \eta=(8N_c)^{-1},
\]
and sufficiently small $\gamma>0$, with the controller budget
$q\le P/(64N_c^2)$. For cut expansion fix $q$ and choose a sufficiently small
positive cut-density constant. The exceptional count covers all clusters,
signs and charts simultaneously, uniformly over $b\ge1$.
\end{theorem}
\begin{outline}
The checked \lean{native_cluster_slab_geometry} applies the bounded-speed
slab selection to the actual append times and native head walks. The carrier
is an integer walk of speed at most $m$. A displacement of size
$v$ crosses slab walls for at most $v$ offsets, so offset averaging bounds
the number of changes by $m(T-1)/W<N_c$. The published short/long slab
dichotomy depends on this change count, not unit speed of the carrier, and
the raw-head stripe estimate is unchanged.

Preserve all protected cores between selected runs and refresh all fringes
at their entries. In the short case all cores start blank. In the long case
pay at most $a(W+m)$ symbols' bits at the outer start. The same calculation
as \cref{eq:slab-margin} gives saving at least
$P/(4N_c)-3N_cdb-a(k+1)$. Each run additionally pays at most
$ak\max(4j+2,2H+1)$ bits for clustered fringes and noncluster windows,
plus its controller. The native replay induction preserves every core together with the selected
control states; matching paid records then recover each demanded output bit.
The typed noncluster imports inject into original completed complementary
returns. Thus \lean{native_bad_cluster_return_bound} derives the finite
input count from the actual source records. Core ownership and current-version replay also prove the localized
cut variant: a latest core writer after the outer start is in the selected
region, and a fringe version is supplied by its current entry snapshot unless
its last visit lies in the current run. Other cases are those in
\cref{sec:fragments}. The checked
\lean{native_bad_cluster_cut_bound} charges the complete local cut record
and omits the import fields entirely; no small-return premise is needed for
this cut variant.

There is no description of the real functions. A visited carrier slab has
$|\ell|\le kT_*+W+1$. At a selected time, chart endpoints differ from the
corresponding oriented head position by at most $j+W$, so their integer
roundings lie within $B_*=kT_*+W+2j+2$. A safe base geometry menu is
\[
 3^k(N_c+1)(2kT_*+W+2)(2B_*+1)^{4k}(T_*+1)^2
 [(T_*+1)^4(2T_*+1)^{k+3}]^{N_c}.
\]
The $3^k$ choices encode exclusion or either orientation of each head.
Consequently its logarithm is $O(\log P)$ at the chosen scale, despite
input-dependent charts. The complete finite exponent is
\[
 P-P/(4N_c)+3N_cdb+a(k+1)
 +aN_ck\max(4j+2,2H+1)+N_cq
 +\log_2\mathcal M_c+\log_2\mathcal S(kT_*,J_0,A).
\]
Put $j=\lceil\varepsilon P\rceil$, $H=\lfloor\varepsilon P\rfloor$,
$J_0=\lfloor\gamma P\rfloor$. Windows/fringes and imports each cost at
most $P/(32N_c)$, apart from constants, and the remaining terms cost at most
$P/(32N_c)$ at large $d$ under the stated state budget. For cuts replace the
sparse-import term by the menu in \cref{eq:local-menu} with $N_c$ in place
of $N$ and choose its cost below $P/(32N_c)$. The checked \lean{cluster_return_rates} and \lean{cluster_cut_rates}
choose one dimension threshold before all positive widths; the resulting
finite exponent is strictly below $(1-\eta)P$. The demand upper bound
uses injectivity of actual append times, and the lower bound uses the
snapshot-adjusted geometric margin. The expansion endpoint covers every
valid actual chart slab and all $H'\le H$ simultaneously, with cuts in the
full original execution graph, including target lower endpoints outside the
selected chunk. An all-head cluster has empty complement and is excluded
on the same return good-input event. The checked singleton chart specializes
the event to every free head; \lean{native_cluster_two_return_heads} and
\lean{native_cluster_one_head_bad} recover the two-head and one-head
consequences of \cref{thm:returns} at these cluster constants.
\end{outline}

\begin{lemma}[Finite order-conflict chart]\label{lem:order-chart}
\formal{OrderChart.lean}{finite_order_chart,order_conflict_le_two_error}
For $m\ge1$ and a nonempty finite family of vectors $v(t)\in\R^m$, put $z(t)=\sum_hv_h(t)$ and
$D=\max_{z(t)\le z(u),h}(v_h(t)-v_h(u))_+$. There is a normalized chart
with $-(m-1)D\le f_h(z(t))-v_h(t)\le D$. Conversely a chart of error $J$
implies $D\le2J$.
\end{lemma}
\begin{outline}
Take coordinatewise prefix maxima over all vectors of sum at most $z(t)$.
These vectors $g(t)$ form a coordinatewise chain, even when sums tie, and
$v_h(t)\le g_h(t)\le v_h(t)+D$. Put $\Gamma_t=\sum_hg_h(t)$.
The formal construction interpolates this chain by a continuous monotone
max-envelope curve. Its $h$th coordinate is
\[
 G_h(x)=\max_t\bigl[g_h(t)-(\Gamma_t-x)_+\bigr]
       +\mathbf 1_{h=h_0}(x-\max_t\Gamma_t)_+,
\]
where $h_0$ is any fixed coordinate. It passes through $g(t)$ at
$x=\Gamma_t$, and its coordinate sum is surjective onto $\R$.
Reparametrize by that sum. Monotonicity makes the reparametrization
well-defined even on constant intervals, and the normalized coordinate sum
makes every coordinate $1$-Lipschitz. Since $z(t)\le\Gamma_t$, the error at
$z(t)$ is at most $D$ in each coordinate. The errors sum to zero, giving the
lower bound. Conversely monotonicity and two error allowances bound every
ordered conflict by $2J$.
\end{outline}
For two work heads, \cref{thm:clusters,lem:order-chart} require macroscopic
same-direction and macroscopic opposite-direction displacements, possibly
on different time intervals: use orientations $(+,-)$ and $(+,+)$. They do
not assert that every two-head path has a small-conflict cluster.
